\chapter{Conclusão}\label{chp:conclusao}
\section{Ameaça à Validade}\label{sec:validation_threat}

% \subsubsection{Tempo Médio Real de Sessão}
% AMEAÇA À VALIDADE: Presença Técnica vs Cognitiva
% A métrica assume que a extensão ativa corresponde a atenção cognitiva do aluno,
% porém casos de "presença fantasma" (aba aberta sem interação real) podem
% superestimar o engajamento. Esta limitação será discutida na Seção de Ameaças à Validade.

% AMEAÇA À VALIDADE: Frequência de Traces  
% O cálculo depende da frequência e consistência dos traces emitidos pela extensão.
% Gaps prolongados entre interações podem indicar problemas técnicos ou desengajamento
% não capturado pela métrica atual.

% \subsubsection{Taxa de Presença}
% AMEAÇA À VALIDADE: Normalização Artificial
% O truncamento em 100% pode mascarar casos onde alunos permanecem além do horário,
% perdendo informação sobre engajamento extra. Alternativas seriam manter a razão
% original ou usar escalas logarítmicas para melhor distribuição.

% AMEAÇA À VALIDADE: Comparação Temporal Assimétrica
% A comparação entre tempo médio individual e duração total da aula pode distorcer
% a percepção de engajamento coletivo, especialmente em turmas com alta variância
% de participação entre alunos.

% \subsubsection{Taxa de Distração}
% AMEAÇA À VALIDADE: Classificação Binária de Distração
% A função is_distraction() opera com classificação binária (distração/não-distração),
% podendo não capturar nuances em sites que podem ser tanto educacionais quanto
% distractores dependendo do contexto de uso.

% AMEAÇA À VALIDADE: Atribuição Temporal Proximal
% O algoritmo atribui todo o intervalo entre dois traces ao estado do trace anterior,
% o que pode superestimar tempos de distração se a mudança de estado ocorrer
% no meio do intervalo.

% AMEAÇA À VALIDADE: Sensibilidade à Frequência de Traces
% Métrica altamente dependente da frequência de emissão de traces - baixa frequência
% pode levar a estimativas imprecisas dos períodos de distração.

\section{Trabalhos Futuros}\label{sec:future_endeavor}
Pretende-se, futuramente, implementar \textit{scaling} vertical na infraestrutura do ESPEON, de modo a suportar turmas maiores e permitir a obtenção de amostras mais representativas. Há também interesse em publicar a extensão na \textit{Chrome Web Store}, possibilitando que docentes de todo o país utilizem a ferramenta de forma gratuita e descomplicada.

Além disso, considera-se a adoção de tecnologias para aprimoramento da interface dos \textit{pop-ups}, especialmente devido ao modelo atual de gerenciamento de estado das telas, realizado por referência no cliente. Por fim, vislumbra-se a criação de \textit{Bases do Conhecimento} para ingestão em \textit{LLMs}, permitindo fornecer aos docentes \textit{insights} sobre suas aulas por meio de diagnósticos em linguagem natural, contribuindo para processos de reflexão e aprimoramento contínuo no aprendizado.
